\documentclass[12pt]{article}
\usepackage[utf8]{inputenc}
\usepackage[T1]{fontenc}
\usepackage[margin=1in]{geometry}
\usepackage{setspace}
\usepackage{xcolor}
\usepackage[colorlinks=true, urlcolor=blue!60!black, linkcolor=blue!60!black]{hyperref}
\setlength{\parindent}{2em}
\setlength{\parskip}{1.25em}
\renewcommand{\baselinestretch}{1.0}
\begin{document}
\begin{center}
{\LARGE \textit{Lemon v.\ Kurtzman}, 403 U.S.\ 602 (1971)\par}
\vspace{0.8em}
{\large Case Brief --- Four-Page Edition\par}
\end{center}
\vspace{1em}

\section{Memory Jogger}

The Court struck down Pennsylvania and Rhode Island programs funding secular instruction in parochial schools because policing the subsidies created excessive entanglement with religion, and announced the three-prong \textit{Lemon} test (secular purpose, primary effect, no excessive entanglement), which governed Establishment Clause doctrine until it was abandoned in 2022 (\href{https://supreme.justia.com/cases/federal/us/403/602/}{Justia}).

\section{Detailed Case Facts}

Two statutes sought to relieve financial pressure on nonpublic schools, most of them Roman Catholic. \textbf{Rhode Island's 1969 Salary Supplement Act} paid up to 15\% supplements directly to certified teachers of secular subjects in nonpublic elementary schools, who had to use only public-school materials and promise in writing not to teach religion; all 250 applicants taught in Catholic schools governed by the Bishop of Providence. \textbf{Pennsylvania's 1968 Act} had the state ``purchase'' secular services by reimbursing nonpublic schools for teachers' salaries, textbooks, and materials in mathematics, modern foreign languages, physical science, and physical education, with state-approved materials and state-audited accounts separating secular from religious costs. It spent about \$5 million a year on 1,181 schools, more than 96\% of whose students attended church-related schools.

\section{Procedural History}

In Rhode Island, a three-judge district court held the Act unconstitutional for fostering ``excessive entanglement.'' In Pennsylvania, a three-judge court recognized taxpayer standing for Alton J.\ Lemon and others but dismissed the complaint over Chief Judge Hastie's dissent. The Supreme Court consolidated the appeals and decided them on June 28, 1971. (A later dispute over reimbursement for past services produced \textit{Lemon II} in 1973.)

\section{Judicial Votes}

\begin{itemize}
\item \textbf{Pennsylvania (No.\ 89):} 8--0, reversed and remanded (Marshall not participating). \textbf{Rhode Island (Nos.\ 569, 570):} 8--1, affirmed.
\item \textbf{Majority:} Chief Justice Burger, joined by Justices Black, Douglas, Harlan, Stewart, Blackmun, and (in the Rhode Island cases) Marshall.
\item \textbf{Concurrences:} Justice Douglas, joined by Black (and Marshall in part); Justice Brennan.
\item \textbf{Concurring in part and dissenting in part:} Justice White (concurring in the judgment in Pennsylvania, dissenting in Rhode Island).
\end{itemize}

\section{Holding}

Both statutes violate the Establishment Clause. A statute must have a secular legislative purpose; its principal or primary effect must neither advance nor inhibit religion; and it must not foster ``an excessive government entanglement with religion,'' judged by the character of the benefited institutions, the nature of the aid, and the resulting relationship between government and religious authority. The decision rested only on entanglement: the Court accepted the secular purposes and did not decide the effects prong. Later cases eroded the test (\textit{Zelman}, 2002), and \textit{Kennedy v.\ Bremerton} (2022) said the Court had ``long ago abandoned \textit{Lemon}.''

\section{Analysis of Opinions}

\textbf{Majority (Burger).} The Religion Clauses are ``at best, opaque,'' but guard against ``sponsorship, financial support, and active involvement of the sovereign in religious activity.'' Unlike bus rides (\textit{Everson}) or secular textbooks (\textit{Allen}), a teacher in a pervasively religious school works under ecclesiastical authority and cannot easily separate faith from secular teaching. Enforcing the bans on religious instruction would require continuing surveillance, compounded by audits of school finances; Pennsylvania's direct payments to church schools were worse still. Annual appropriations benefiting mainly one faith would also foster ``political division along religious lines.''

\textbf{Douglas (concurring).} Drawing on Madison's \textit{Memorial and Remonstrance}, Douglas called parochial schools ``an organism living on one budget'': subsidies free church funds for religious teaching, while policing them threatens both Religion Clauses.

\textbf{Brennan (concurring).} Direct subsidies to religious schools differ from tax exemptions or student benefits; the statutes used ``essentially religious means to serve governmental ends, where secular means would suffice.''

\textbf{White (dissenting in part).} States may fund the separable secular function of parochial schools; the Rhode Island court found no religion injected into secular classes. The Court created an ``insoluble paradox'' by condemning both unsafeguarded aid and the safeguards themselves, and could not square its result with \textit{Tilton}.

\section{Examples: Future Applications}

\begin{enumerate}
\item \textbf{Same-Side: STEM Salary Subsidies.} A state pays part of religious-school science teachers' salaries, with classroom inspections and budget audits. Direct funding plus surveillance is excessive entanglement.
\item \textbf{Same-Side: Classroom Creeds.} A law requiring a state-selected religious creed in every classroom, adopted to promote a faith, fails the purpose and effect prongs.
\item \textbf{Opposite-Side: Neutral Vouchers.} Scholarships usable at any school by genuine private choice, with no monitoring of content, are valid under \textit{Zelman}.
\item \textbf{Opposite-Side: General Services.} Fire inspections and bus transportation for all schools are permitted contacts under \textit{Everson} and \textit{Walz}.
\item \textbf{Fence-Sitter: Public Tutors on Parochial Campuses.} Remedial specialists in rooms stripped of symbols were first struck down under \textit{Lemon}'s presumption, but \textit{Agostini v.\ Felton} (1997) rejected that presumption.
\end{enumerate}

\section{Critique}

\textbf{Originalist critique.} Critics call the test ahistorical and manipulable. Justice Scalia likened it to ``some ghoul in a late-night horror movie that repeatedly sits up in its grave and shuffles abroad, after being repeatedly killed and buried,'' and critics echoing White argued that it penalized religious families, culminating in \textit{Kennedy v.\ Bremerton}'s turn to ``historical practices and understandings.''

\textbf{Separationist defense.} Progressive scholars defend \textit{Lemon}'s insight that barring sponsorship and entanglement protects religious minorities and religious institutions, and warn that an undefined history test privileges majority faiths.

\section{Key Quotations}

\begin{quote}\sloppy
``Every analysis in this area must begin with consideration of the cumulative criteria developed by the Court over many years. Three such tests may be gleaned from our cases. First, the statute must have a secular legislative purpose; second, its principal or primary effect must be one that neither advances nor inhibits religion, \textit{Board of Education v.\ Allen}, 392 U.S.\ 236, 243 (1968); finally, the statute must not foster `an excessive government entanglement with religion.' \textit{Walz}, \textit{supra}, at 674.''
\end{quote}
\noindent --- Chief Justice Burger, for the Court.

\begin{quote}
``Judicial caveats against entanglement must recognize that the line of separation, far from being a `wall,' is a blurred, indistinct, and variable barrier depending on all the circumstances of a particular relationship.''
\end{quote}
\noindent --- Chief Justice Burger, for the Court.

\begin{quote}\sloppy
``A comprehensive, discriminating, and continuing state surveillance will inevitably be required to ensure that these restrictions are obeyed and the First Amendment otherwise respected. Unlike a book, a teacher cannot be inspected once so as to determine the extent and intent of his or her personal beliefs and subjective acceptance of the limitations imposed by the First Amendment.''
\end{quote}
\noindent --- Chief Justice Burger, for the Court.

\end{document}
